\documentclass[10pt,a4paper,twoside,bg=print,twocolumn,openany,nodeprecatedcode]{dndbook}
\usepackage[utf8]{inputenc}
\usepackage{tabulary}
\usepackage[normalem]{ulem}
\usepackage{color,soul}
\usepackage{float}
\usepackage{caption}
\usepackage{wrapfig}
\usepackage{tocloft}
\usepackage[toc]{multitoc}
\usepackage[hidelinks]{hyperref}
\raggedbottom
\extrafloats{2000}
\cftsetindents{section}{0em}{0em}
\cftsetindents{subsection}{0em}{0em}
\cftsetindents{subsubsection}{0.2em}{0em}
\setcounter{tocdepth}{1}
\renewcommand*{\multicolumntoc}{3}
\setlist[description]{nolistsep,listparindent=3pt,topsep=6pt}
\date{}
\begin{document}
\frontmatter
\tableofcontents
\mainmatter
\chapter{Using This Book}
\DndDropCapLine{T}{}he multiverse of Dungeons and Dragons brims with peril, and many of the greatest dangers are monsters. A companion to the \textit{Monster Manual}, this book collects monsters from many different planes of existence, each creature ready to imperil D\&D heroes of different levels. The book also includes game statistics for nonplayer characters, who may assist or oppose the heroes, and an array of fantastical races for players to consider for their characters. The monsters and fantastical races in the following chapters are accompanied by the occasional comments of the mighty wizard Mordenkainen, one of the most learned mages of the D\&D world of Greyhawk. The archmage Tasha, Mordenkainen's friendly rival, interjects as well. All together, the monsters, nonplayer characters, and fantastical races herein provide a host of new friends and foes to populate your D\&D worlds.


\section{What You'll Find Within}

\begin{description}
\item[Chapter 1, "Fantastical Races,".]
presents over 30 race options for player characters, complementing the options in the \textit{Player's Handbook} and other D&D books. These races debuted elsewhere and appear all together for the first time here, each of them revised to fit into the current state of the game.
\item[Chapter 2, "Bestiary,".]
contains over 250 monsters and NPCs, each one represented by a stat block and story text. When you're preparing to run an adventure as the DM, consider sprinkling these creatures into your games, mixing them with the monsters and NPCs from the \textit{Monster Manual}. The creatures in this chapter fit seamlessly with the ones you use from other D&D books.
\item[Appendix, "Monster Lists,".]
makes it easy for you to find the right stat blocks for your adventures. This appendix lists the book's monsters and NPCs by creature type, challenge rating, and environment.
\end{description}

\chapter{Fantastical Races}
\DndDropCapLine{G}{}athering together fantastical races from throughout the D\&D multiverse, this chapter offers the following races for player characters, supplementing the race options in the \textit{Player's Handbook}:



\begin{itemize}
\begin{multicols}{4}
\item Aarakocra
\item Aasimar
\item Bugbear
\item Centaur
\item Changeling
\item Deep Gnome
\item Duergar
\item Eladrin
\item Fairy
\item Firbolg
\item Genasi (Air)
\item Genasi (Earth)
\item Genasi (Fire)
\item Genasi (Water)
\item Githyanki
\item Githzerai
\item Goblin
\item Goliath
\item Harengon
\item Hobgoblin
\item Kenku
\item Kobold
\item Lizardfolk
\item Minotaur
\item Orc
\item Satyr
\item Sea Elf
\item Shadar-kai
\item Shifter
\item Tabaxi
\item Tortle
\item Triton
\item Yuan-ti
\end{multicols}

\end{itemize}

Many of these races are based on creatures that appear in the \textit{Monster Manual} or the bestiary of this book. Consult with your DM to see whether an option here is appropriate for your campaign. If you do use a race in this chapter, first read the "Creating Your Character" section below.

\section{Creating Your Character}
At 1st level, you choose whether your character is a member of the human race or of a fantastical race. If you select a fantastical race in this chapter, follow these additional rules during character creation.

\subsection{Ability Score Increases}
When determining your character's ability scores, increase one score by 2 and increase a different score by 1, or increase three different scores by 1. Follow this rule regardless of the method you use to determine the scores, such as rolling or point buy. The "Quick Build" section for your character's class offers suggestions on which scores to increase. You can follow those suggestions or ignore them, but you can't raise any of your scores above 20.

\subsection{Languages}
Your character can speak, read, and write Common and one other language that you and your DM agree is appropriate for the character. The Player's Handbook offers a list of languages to choose from. The DM is free to modify that list for a campaign.

\subsection{Creature Type}
Every creature in D\&D, including each player character, has a special tag in the rules that identifies the type of creature they are. Most player characters are of the Humanoid type. A race in this chapter tells you what your character's creature type is. Here's a list of the game's creature types in alphabetical order: Aberration, Beast, Celestial, Construct, Dragon, Elemental, Fey, Fiend, Giant, Humanoid, Monstrosity, Ooze, Plant, Undead. These types don't have rules themselves, but some rules in the game affect creatures of certain types in different ways. For example, the \textit{cure wounds} spell doesn't work on a Construct or an Undead.

\subsection{Life Span}
The typical life span of a player character in the D\&D multiverse is about a century, assuming the character doesn't meet a violent end on an adventure. Members of some races, such as dwarves and elves, can live for centuries. If typical members of a race in this book can live longer than a century, that fact is mentioned in the race's description.

\subsection{Height and Weight}
Player characters, regardless of race, typically fall into the same ranges of height and weight that humans have in our world. If you'd like to determine your character's height or weight randomly, consult the Random Height and Weight table in the Player's Handbook, and choose the row in the table that best represents the build you imagine for your character.

\chapter{Bestiary}
\DndDropCapLine{T}{}his bestiary provides game statistics and lore for more than 250 monsters, which are suitable for any D\&D campaign. The chapter includes old favorites from past editions of the game as well as creatures created for the current edition.


This chapter is a companion to the \textit{Monster Manual} and adopts a similar presentation. If you are unfamiliar with the monster stat block format, read the introduction of the \textit{Monster Manual} before proceeding further. It explains stat block terminology and gives rules for various monster traits—information that isn't repeated here.

The creatures in this bestiary are organized alphabetically. A few are gathered under a group heading; for example, the "Dinosaurs" section contains stat blocks for various dinosaurs.

Herein you will find some weapons that deal unusual damage types and spellcasting that functions in atypical ways. Such an exception is a special feature of a monster and represents how it uses the weapon or casts its spells; the exception has no affect on how a weapon or a spell functions for someone else.

Finally, if a stat block contains the name of a class in the monster's name or in parentheses under the name, the monster is considered a member of that class for the purpose of meeting prerequisites for magic items.

\DndQuote%
{}
{''My friend Mordenkainen constantly fights monsters and makes notes about them. I love notes as much as the next wizard, but fighting gets old. Why not sit down and have a drink with the monster? I've learned the most delicious secrets from my monstrous drinking companions. The other monsters? I disintegrated them.''}
{Tasha}

\begin{itemize}
\begin{multicols}{4}
\item \textbf{Black Abishai}
\item \textbf{Blue Abishai}
\item \textbf{Green Abishai}
\item \textbf{Red Abishai}
\item \textbf{White Abishai}
\item \textbf{Alhoon}
\item \textbf{Alkilith}
\item \textbf{Allip}
\item \textbf{Amnizu}
\item \textbf{Annis Hag}
\item \textbf{Archdruid}
\item \textbf{Archer}
\item \textbf{Armanite}
\item \textbf{Astral Dreadnought}
\item \textbf{Babau}
\item \textbf{Bael}
\item \textbf{Balhannoth}
\item \textbf{Banderhobb}
\item \textbf{Baphomet}
\item \textbf{Bard}
\item \textbf{Barghest}
\item \textbf{Berbalang}
\item \textbf{Bheur Hag}
\item \textbf{Blackguard}
\item \textbf{Bodak}
\item \textbf{Boggle}
\item \textbf{Boneclaw}
\item \textbf{Bulezau}
\item \textbf{Cadaver Collector}
\item \textbf{Canoloth}
\item \textbf{Catoblepas}
\item \textbf{Aurochs}
\item \textbf{Deep Rothé}
\item \textbf{Ox}
\item \textbf{Stench Kow}
\item \textbf{Cave Fisher}
\item \textbf{Champion}
\item \textbf{Chitine}
\item \textbf{Choker}
\item \textbf{Choldrith}
\item \textbf{Clockwork Bronze Scout}
\item \textbf{Clockwork Iron Cobra}
\item \textbf{Clockwork Oaken Bolter}
\item \textbf{Clockwork Stone Defender}
\item \textbf{Cloud Giant Smiling One}
\item \textbf{Corpse Flower}
\item \textbf{Cranium Rat}
\item \textbf{Swarm of Cranium Rats}
\item \textbf{Darkling}
\item \textbf{Darkling Elder}
\item \textbf{Death Kiss}
\item \textbf{Deathlock}
\item \textbf{Deathlock Mastermind}
\item \textbf{Deathlock Wight}
\item \textbf{Deep Scion}
\item \textbf{Demogorgon}
\item \textbf{Derro}
\item \textbf{Derro Savant}
\item \textbf{Devourer}
\item \textbf{Dhergoloth}
\item \textbf{Brontosaurus}
\item \textbf{Deinonychus}
\item \textbf{Dimetrodon}
\item \textbf{Hadrosaurus}
\item \textbf{Quetzalcoatlus}
\item \textbf{Stegosaurus}
\item \textbf{Velociraptor}
\item \textbf{Dolphin}
\item \textbf{Dolphin Delighter}
\item \textbf{Draegloth}
\item \textbf{Drow Arachnomancer}
\item \textbf{Drow Favored Consort}
\item \textbf{Drow House Captain}
\item \textbf{Drow Inquisitor}
\item \textbf{Drow Matron Mother}
\item \textbf{Drow Shadowblade}
\item \textbf{Duergar Despot}
\item \textbf{Duergar Kavalrachni}
\item \textbf{Duergar Mind Master}
\item \textbf{Duergar Soulblade}
\item \textbf{Duergar Stone Guard}
\item \textbf{Duergar Screamer}
\item \textbf{Duergar Warlord}
\item \textbf{Duergar Xarrorn}
\item \textbf{Duergar Hammerer}
\item \textbf{Dybbuk}
\item \textbf{Eidolon}
\item \textbf{Sacred Statue}
\item \textbf{Autumn Eladrin}
\item \textbf{Spring Eladrin}
\item \textbf{Summer Eladrin}
\item \textbf{Winter Eladrin}
\item \textbf{Elder Brain}
\item \textbf{Elder Tempest}
\item \textbf{Air Elemental Myrmidon}
\item \textbf{Earth Elemental Myrmidon}
\item \textbf{Fire Elemental Myrmidon}
\item \textbf{Water Elemental Myrmidon}
\item \textbf{Fire Giant Dreadnought}
\item \textbf{Firenewt Warlock of Imix}
\item \textbf{Firenewt Warrior}
\item \textbf{Flail Snail}
\item \textbf{Flind}
\item \textbf{Fraz-Urb'luu}
\item \textbf{Froghemoth}
\item \textbf{Frost Giant Everlasting One}
\item \textbf{Frost Salamander}
\item \textbf{Gauth}
\item \textbf{Gazer}
\item \textbf{Geryon}
\item \textbf{Giant Strider}
\item \textbf{Giff}
\item \textbf{Girallon}
\item \textbf{Githyanki Gish}
\item \textbf{Githyanki Kith'rak}
\item \textbf{Githyanki Supreme Commander}
\item \textbf{Githzerai Anarch}
\item \textbf{Githzerai Enlightened}
\item \textbf{Gnoll Flesh Gnawer}
\item \textbf{Gnoll Hunter}
\item \textbf{Gnoll Witherling}
\item \textbf{Gray Render}
\item \textbf{Graz'zt}
\item \textbf{Grung}
\item \textbf{Grung Elite Warrior}
\item \textbf{Grung Wildling}
\item \textbf{Guard Drake}
\item \textbf{Hellfire Engine}
\item \textbf{Hobgoblin Devastator}
\item \textbf{Hobgoblin Iron Shadow}
\item \textbf{Howler}
\item \textbf{Hutijin}
\item \textbf{Hydroloth}
\item \textbf{Juiblex}
\item \textbf{Ki-rin}
\item \textbf{Kobold Dragonshield}
\item \textbf{Kobold Inventor}
\item \textbf{Kobold Scale Sorcerer}
\item \textbf{Korred}
\item \textbf{Kraken Priest}
\item \textbf{Young Kruthik}
\item \textbf{Adult Kruthik}
\item \textbf{Kruthik Hive Lord}
\item \textbf{Leucrotta}
\item \textbf{Leviathan}
\item \textbf{Martial Arts Adept}
\item \textbf{Marut}
\item \textbf{Master Thief}
\item \textbf{Maurezhi}
\item \textbf{Maw Demon}
\item \textbf{Meazel}
\item \textbf{Meenlock}
\item \textbf{Merregon}
\item \textbf{Merrenoloth}
\item \textbf{Mindwitness}
\item \textbf{Moloch}
\item \textbf{Molydeus}
\item \textbf{Morkoth}
\item \textbf{Mouth of Grolantor}
\item \textbf{Nabassu}
\item \textbf{Nagpa}
\item \textbf{Narzugon}
\item \textbf{Neogi Hatchling}
\item \textbf{Neogi}
\item \textbf{Neogi Master}
\item \textbf{Neothelid}
\item \textbf{Nightwalker}
\item \textbf{Nilbog}
\item \textbf{Nupperibo}
\item \textbf{Oblex Spawn}
\item \textbf{Adult Oblex}
\item \textbf{Elder Oblex}
\item \textbf{Ogre Battering Ram}
\item \textbf{Ogre Bolt Launcher}
\item \textbf{Ogre Chain Brute}
\item \textbf{Ogre Howdah}
\item \textbf{Oinoloth}
\item \textbf{Orcus}
\item \textbf{Orthon}
\item \textbf{Phoenix}
\item \textbf{Quickling}
\item \textbf{Redcap}
\item \textbf{Retriever}
\item \textbf{Rutterkin}
\item \textbf{Sea Spawn}
\item \textbf{Shadar-kai Gloom Weaver}
\item \textbf{Shadar-kai Shadow Dancer}
\item \textbf{Shadar-kai Soul Monger}
\item \textbf{Shadow Mastiff}
\item \textbf{Shadow Mastiff Alpha}
\item \textbf{Shoosuva}
\item \textbf{Sibriex}
\item \textbf{Skulk}
\item \textbf{Skull Lord}
\item \textbf{Slithering Tracker}
\item \textbf{Angry Sorrowsworn}
\item \textbf{Hungry Sorrowsworn}
\item \textbf{Lonely Sorrowsworn}
\item \textbf{Lost Sorrowsworn}
\item \textbf{Wretched Sorrowsworn}
\item \textbf{Spawn of Kyuss}
\item \textbf{Star Spawn Grue}
\item \textbf{Star Spawn Hulk}
\item \textbf{Star Spawn Larva Mage}
\item \textbf{Star Spawn Mangler}
\item \textbf{Star Spawn Seer}
\item \textbf{Female Steeder}
\item \textbf{Male Steeder}
\item \textbf{Steel Predator}
\item \textbf{Stone Cursed}
\item \textbf{Stone Giant Dreamwalker}
\item \textbf{Storm Giant Quintessent}
\item \textbf{Swarm of Rot Grubs}
\item \textbf{Swashbuckler}
\item \textbf{Sword Wraith Commander}
\item \textbf{Sword Wraith Warrior}
\item \textbf{Tanarukk}
\item \textbf{Titivilus}
\item \textbf{Tlincalli}
\item \textbf{Tortle}
\item \textbf{Tortle Druid}
\item \textbf{Trapper}
\item \textbf{Dire Troll}
\item \textbf{Rot Troll}
\item \textbf{Spirit Troll}
\item \textbf{Venom Troll}
\item \textbf{Ulitharid}
\item \textbf{Vampiric Mist}
\item \textbf{Vargouille}
\item \textbf{Vegepygmy}
\item \textbf{Thorny Vegepygmy}
\item \textbf{Vegepygmy Chief}
\item \textbf{War Priest}
\item \textbf{Warlock of the Archfey}
\item \textbf{Warlock of the Fiend}
\item \textbf{Warlock of the Great Old One}
\item \textbf{Warlord}
\item \textbf{Wastrilith}
\item \textbf{Apprentice Wizard}
\item \textbf{Abjurer Wizard}
\item \textbf{Conjurer Wizard}
\item \textbf{Diviner Wizard}
\item \textbf{Enchanter Wizard}
\item \textbf{Evoker Wizard}
\item \textbf{Illusionist Wizard}
\item \textbf{Necromancer Wizard}
\item \textbf{Transmuter Wizard}
\item \textbf{Wood Woad}
\item \textbf{Xvart}
\item \textbf{Xvart Warlock of Raxivort}
\item \textbf{Yagnoloth}
\item \textbf{Yeenoghu}
\item \textbf{Yeth Hound}
\item \textbf{Yuan-ti Anathema}
\item \textbf{Yuan-ti Broodguard}
\item \textbf{Yuan-ti Mind Whisperer}
\item \textbf{Yuan-ti Nightmare Speaker}
\item \textbf{Yuan-ti Pit Master}
\item \textbf{Zaratan}
\item \textbf{Zariel}
\item \textbf{Zuggtmoy}
\end{multicols}

\end{itemize}

\chapter{Maps}
\chapter{Credits}

\begin{description}
\begin{multicols}{2}

\begin{description}
\item[Lead Designer.]
\DndDropCapLine{J}{}eremy Crawford


\item[Art Directors.]
Emi Tanji, Kate Irwin

\item[Original Design.]
Jeremy Crawford, Adam Lee, Mike Mearls, Christopher Perkins, Ben Petrisor, Sean K Reynolds, Robert J. Schwalb, Matt Sernett, Chris Sims, Nolan Whale, Steve Winter

\item[Revision Design.]
Sydney Adams, Judy Bauer, Jeremy Crawford, Makenzie De Armas, Dan Dillon, Ari Levitch, Ben Petrisor, Taymoor Rehman

\item[Editing.]
Judy Bauer, Michele Carter, Kim Mohan, Christopher Perkins

\item[Graphic Designers.]
Trystan Falcone, Emi Tanji

\item[Cover Illustrator.]
Grzegorz Rutkowski

\item[Interior Illustrators.]
Dave Allsop, Tom Babbey, John-Paul Balmet, Thomas M. Baxa, Mark Behm, Eric Belisle, Michael Berube, Mike Bierek, Jared Blando, Zoltan Boros, Christopher Bradley, Alix Branwyn, Aleksi Briclot, IP Dmitry Burmak, Filip Burburan, Christopher Burdett, Sidharth Chaturvedi, Jedd Chevrier, Conceptopolis, Daarken, Nikki Dawes, Eric Deschamps, Simon Dominic, Dave Dorman, Nicholas Elias, Wayne England, Jason Felix, Justin Gerard, Justyna Gil, Lars Grant-West, Jon Hodgson, Ralph Horsley, Tyler Jacobson, Mike Jordana, Vance Kelly, Julian Kok, Michael Komarck, Daniel Landerman, Olly Lawson, Daniel Ljunggren, Valera Lutfullina, Warren Mahy, Lorenzo Mastroianni, Brynn Metheney, Aaron Miller, Francisco Miyara, Caio Monteiro, Scott Murphy, Marta Nael, Marco Nelor, Jim Nelson, Adam Paquette, PINDURSKI, Claudio Pozas, April Prime, Grzegorz Rutkowski, Marc Sasso, Chris Seaman, Ilya Shkipin, Rudy Siswanto, David Sladek, Craig J Spearing, Bryan Sola, Zack Stella, Sarah Stone, Philip Straub, Arnie Swekel, Thom Tenery, David A. Trampier, Cory Trego-Erdner, Beth Trott, Brian Valeza, Randy Vargas, Franz Vohwinkel, Anthony S. Waters, Campbell White, Richard Whitters, Eva Widermann, Sam Wood, Shawn Wood, Ben Wootten, Zuzanna Wuuyk, Min Yum

\item[Project Engineer.]
Cynda Callaway

\item[Imaging Technician.]
Kevin Yee

\item[Prepress Specialist.]
Jefferson Dunlap

\end{description}

\section{D\&D Studio}

\begin{description}
\item[Executive Producer.]
Ray Winninger

\item[Principal Designers.]
Jeremy Crawford, Christopher Perkins

\item[Design Manager.]
Steve Scott

\item[Design Department.]
Sydney Adams, Judy Bauer, Makenzie De Armas, Dan Dillon, Amanda Hamon, Ari Levitch, Ben Petrisor, Taymoor Rehman, F. Wesley Schneider, James Wyatt

\item[Senior Art Director.]
Richard Whitters

\item[Art Department.]
Trystan Falcone, Kate Irwin, Emi Tanji, Shawn Wood, Trish Yochum

\item[Senior Producer.]
Dan Tovar

\item[Producers.]
Bill Benham, Robert Hawkey, Lea Heleotis

\item[Director of Product Management.]
Liz Schuh

\item[Product Managers.]
Natalie Egan, Chris Lindsay, Hilary Ross

\end{description}

\section{D\&D Marketing}

\begin{description}
\item[Director of Global Brand Marketing.]
Brian Perry

\item[Global Brand Manager.]
Shelly Mazzanoble

\item[Senior Marketing Communications Manager.]
Greg Tito

\item[Community Manager.]
Brandy Camel

\end{description}

\section{D\&D Beyond}

\begin{description}
\item[Digital Design Manager.]
Faith Elisabeth Lilley

\item[Digital Design Team.]
Jay Jani, Adam Walton

\end{description}

\end{multicols}

\end{description}

This book is a revision of content that originally appeared in the books \textit{Volo's Guide to Monsters} (2016) and \textit{Mordenkainen's Tome of Foes} (2018). It also includes revised options from Princes of the Apocalypse (2015), \textit{Eberron: Rising from the Last War} (2019), and \textit{Mythic Odysseys of Theros} (2020).

\end{document}
